\section{Text Preprocessing and Feature Representation}
\label{sec:text-features}

\paragraph{Representation.}
We used the fixed 796-review training partition and reserved the
200-review test partition for subsequent final evaluation.
After the shared data cleaning, we lowercased text and applied Unicode
accent stripping. The word tokenizer retained tokens of at least two
word characters. We did not remove stop words or apply stemming or
lemmatization, retaining explicit negation words such as \emph{not},
\emph{no}, and \emph{never}. The default tokenizer nevertheless splits
contractions and drops single-character fragments; this does not
preserve all contracted negation. Punctuation intensity and long-range
word order are not represented.

We compared word unigrams with word unigrams plus adjacent bigrams using
scikit-learn's TF--IDF implementation \citep{pedregosa2011scikit,manning2008ir}.
Term frequency was scaled as $1+\ln(c)$ for positive count $c$.
Smoothed inverse document frequency was
$\ln((1+n)/(1+\mathrm{df}))+1$, followed by document-wise L2 normalization.
We set \texttt{min\_df=1}, \texttt{max\_df=1.0}, and imposed no vocabulary
cap. The resulting nonnegative features are compatible with both
Multinomial Naive Bayes and Linear SVM.

\paragraph{Training-only comparison.}
We used identical five-fold stratified training splits with shuffling
and random seed 8017. A fresh vectorizer and classifier were fitted
inside every fold; validation texts did not determine vocabulary or IDF.
Linear SVM with fixed $C=1$ was the feature-selection probe, and
Multinomial Naive Bayes with fixed $\alpha=1$ was a sensitivity check.
We selected the representation by mean positive-class validation F1
under the SVM probe, preferring unigrams in an exact tie. This experiment
does not replace later classifier tuning or final test evaluation.

\begin{table}[t]
  \caption{Training-only five-fold feature comparison. Values are mean
  positive-class F1 and sample standard deviation across folds, not
  confidence intervals. No test-set performance is reported.}
  \label{tab:tfidf-ablation}
  \centering
  \resizebox{\columnwidth}{!}{%
  \begin{tabular}{llr}
    \hline
    Features & Probe & F1 mean $\pm$ SD \\
    \hline
    Unigram & Linear SVM & $0.8016 \pm 0.0282$ \\
    Unigram + bigram & Linear SVM & $0.7998 \pm 0.0388$ \\
    Unigram & Multinomial NB & $0.7962 \pm 0.0484$ \\
    Unigram + bigram & Multinomial NB & $0.7845 \pm 0.0394$ \\
    \hline
  \end{tabular}
  }
\end{table}

\paragraph{Feature summary and handoff.}
Unigrams yielded a mean SVM validation F1 of 0.8016, compared with
0.7998 for unigrams plus bigrams (Table~\ref{tab:tfidf-ablation}).
The difference of approximately 0.0018 was small relative to fold
variability and does not establish a significant or general advantage.
Although bigrams slightly increased mean SVM accuracy
(0.8041 versus 0.8015) and ROC-AUC (0.8972 versus 0.8953),
they did not improve the prespecified F1 criterion.
We therefore recommend unigrams as the shared default.

Fitting on the full training partition produced 1,781 unigram features
and 7,345 unigram--bigram features, a 4.12-fold increase in dimension.
The training matrices contained 7,836 and 15,269 nonzero entries,
respectively, with densities of 0.553\% and 0.261\%.
All training and transformed test rows had at least one nonzero feature.
The larger phrase vocabulary may be difficult to estimate reliably
from this small corpus. We exported sparse matrices, ordered row IDs,
vocabularies, fitted vectorizers, and parameter metadata for subsequent
members. Test texts were transformed only after fitting on training
texts. Classifier cross-validation must still refit the full pipeline
within each fold rather than reuse these full-training fitted matrices.
